\section{Lagwise packet--HLZ bridge for the log-only source}
\label{sec:lagwise-packet-hlz}

This section records the proposed lagwise reduction and the exact geometry that
motivated it.  The reduction is **not presently a theorem**: after the
one-$\chi$ stationary step the physical arithmetic carrier retains an additive
phase, and the untwisted scalar HLZ diagonal of
Section~\ref{sec:scalar-hlz-core} does not by itself remove that phase.  The
precise obstruction and a rank-versus-floor no-go test are recorded in
\texttt{doc/lagwise\_hlz\_gap\_audit.md}.  We keep the formulas below because
they identify the candidate endpoint geometry that any repaired arithmetic
bridge would have to reproduce.

For the principal logarithmic source we keep the packet polarization in its
actual continuous lag variable $d=y-x$.  No Fourier multiplier of a lag
convolution is identified with a one-sided packet-prefix multiplier.

Fix the stable logarithmic core
\begin{equation}
\boxed{
W_T:=L-A_0\log L,
\qquad
X_T:=e^{W_T}=\frac{T}{2\pi L^{A_0}}.
}
\label{eq:stable-log-core-main}
\end{equation}
On the translated flat plateau define the packet synthesis
\begin{equation}
\boxed{
(\mathcal A_Tv)(x):=\phi_L(x)
\sum_{\nu\in\mathcal I_T}v_\nu e^{-i\tau_\nu x}.
}
\label{eq:lagwise-packet-analysis}
\end{equation}
For the candidate principal core, both packet legs are first restricted to the
same stable flat interval; every term with an exterior leg belongs to the
geometric leakage ledger of Section~\ref{sec:shell-frame}.

\begin{proposition}[Candidate lagwise packet--HLZ bridge; open]
\label{prop:lagwise-packet-HLZ}
The following statement is the missing analytic bridge required by the current
log-only route.  It is **not proved** by the scalar HLZ lemma as presently
written.

Restrict to the invariant source $h=k=1$, one retained height block, and the
stable flat packet core. Put
\[
X(\tau)=\frac{\tau}{2\pi},\qquad
\lambda(\tau)=\frac{\log X(\tau)}L,
\]
and, for $|d|\le W_T$, define
\begin{equation}
\boxed{\rho_\tau(d):=\lambda(\tau)-\frac{|d|}{L}.}
\label{eq:lagwise-rho}
\end{equation}
The desired conclusion would be that the logarithmic source $H_{\log}$ of
\eqref{eq:Hlog} has principal stationary packet form
\begin{align}
\mathcal Q_{\log,T}(v,w;p)
={}&\iint
 (\mathcal A_Tv)(x)\overline{(\mathcal A_Tw)(y)}
 \mathcal G^{\log}_{\lambda(\tau),\rho_\tau(y-x)}(p)\,dx\,dy
\notag\\
&+\mathcal R^{\rm HLZ}_{\log,T}(v,w;p),
\label{eq:lagwise-HLZ-form}
\end{align}
where
\begin{equation}
\boxed{
\mathcal G^{\log}_{\lambda,\rho}(p)
=e^{\lambda p}\left[
-2E_\rho(p)^2
+e^{2\lambda}E_\rho(p+2)E_\rho(p-2)
\right],
\qquad
E_\rho(q)=\int_0^\rho e^{-qt}\,dt,
}
\label{eq:logonly-local-profile}
\end{equation}
and that the scalar error profile and its first two $\rho$-derivatives have
arbitrary logarithmic saving uniformly in the retained lag and height blocks.
A proof must, in addition to the scalar HLZ contour deformation, control the
additive physical carrier displayed below.
\end{proposition}

\subsection{Exact lag geometry and the missing carrier}

For basis packets $e_r,e_s$, exact polarization gives
\begin{equation}
\mathcal B_{e_r,e_s}(c+it)
=\iint\phi_L(x)\phi_L(y)e^{-i\tau_sx+i\tau_ry}
 e^{\sigma_0(x-y)}e^{it(x-y)}\,dx\,dy,
\label{eq:lagwise-polarization-recall}
\end{equation}
with $\sigma_0=c-1/2$. For one Dirichlet monomial $\xi^{-c-it}$ put
$d=y-x$. Then exactly
\[
e^{\sigma_0(x-y)}e^{it(x-y)}\xi^{-c-it}
=e^{d/2}(\xi e^d)^{-c-it}.
\]
Hence the stationary dilation parameter of Section~\ref{sec:stationary} is
\begin{equation}
\boxed{w=d=y-x.}
\label{eq:true-packet-lag}
\end{equation}
The one-$\chi$ stationary main term is therefore not merely a cutoff at
$t_*=2\pi e^d\xi$; it also contains the oscillatory factor
\begin{equation}
\boxed{e(-e^d\xi).}
\label{eq:lagwise-additive-carrier}
\end{equation}
When the invariant arithmetic diagonal is exposed, $\xi$ contains the physical
integer carrier (in the old common-product notation this is $N=\ell$).  Thus
\eqref{eq:lagwise-additive-carrier} is an arithmetic additive twist, not a
packet-independent scalar.  The untwisted Euler product
\eqref{eq:scalar-HLZ-diagonal} does not contain this factor.

Ignoring this obstruction only for the purpose of recording the candidate
endpoint geometry, Lemma~\ref{lem:scalar-HLZ-local-height} would put the scalar
diagonal on the local scale $X(\tau)$.  The two Hermitian stationary
orientations then have cutoffs
\[
\ell\le X(\tau)e^{-d},\qquad
\ell\le X(\tau)e^{d},
\]
whose intersection is
\begin{equation}
\boxed{\ell\le X(\tau)e^{-|d|}.}
\label{eq:lagwise-common-cutoff}
\end{equation}
Writing $t=\log\ell/L$ gives the candidate interval
$0\le t\le\rho_\tau(d)$.  This geometric intersection remains correct as a
cutoff statement; what is missing is an arithmetic theorem which contracts the
additively twisted carrier on this interval to the untwisted profile
\eqref{eq:logonly-local-profile} with the required operator error.

If such a theorem were available, the remaining scalar Mellin weight would
have the bounded-support form
\begin{equation}
W_d(u,\mathbf z,p)
=\int_0^{\rho_\tau(d)}
 a_{d,\mathbf z,p,T}(t)e^{Lut}\,dt,
\label{eq:lagwise-Mellin-weight}
\end{equation}
or a fixed finite product of such factors, and the stationary-symbol and
separated-Cauchy estimates would give the polylogarithmic derivative bounds
required by Lemma~\ref{lem:scalar-lagweight-HLZ}.  These observations concern
the smooth amplitude only; they do not control
\eqref{eq:lagwise-additive-carrier}.

\begin{corollary}[Conditional frozen log-only lag kernel]
\label{cor:frozen-logonly-lag-kernel}
**Conditional on Proposition~\ref{prop:lagwise-packet-HLZ}.**  Let
\begin{equation}
\boxed{
G_{\log}(\rho):=[p^2]\,
 e^p\left[-2E_\rho(p)^2+e^2E_\rho(p+2)E_\rho(p-2)\right].
}
\label{eq:Glog-from-source}
\end{equation}
After freezing $\lambda(\tau)=1$ and rescaling $r=x/L$, $s=y/L$, the proposed
principal operator is the pullback of the lag kernel
\begin{equation}
\boxed{K_{\log}(r,s)=G_{\log}(1-|r-s|).}
\label{eq:Klog-lag-kernel}
\end{equation}
The finite-height profile defect would have first two normalized lag
derivatives $O(L^{-1})$.
\end{corollary}

\begin{proof}
Assuming the missing bridge, the map
$(\lambda,\rho,p)\mapsto\mathcal G^{\log}_{\lambda,\rho}(p)$ is analytic on a
fixed neighbourhood of the retained parameter set and
$\lambda(\tau)=1+O(L^{-1})$.  The mean-value theorem gives the parameter defect,
and substitution of $d=y-x$ gives the displayed conditional lag kernel.
\end{proof}

\begin{remark}[What a repair must prove]
A repair cannot consist only of replacing the old one-sided packet multiplier
by a continuous lag variable.  It must keep the carrier
$e(-e^d\ell)$ and show, by an additively twisted Estermann/Voronoi analysis, a
second stationary phase, or another exact mechanism, how the resulting
operator produces the claimed bulk profile with an error compatible with the
relative Hilbert--Schmidt ledger.  Until that step is supplied, the conditional
corollary above and every downstream use of $P_T=A_T^{G_{\log}}$ do not prove
the main theorem.
\end{remark}
